\chapter{Belajar dari Data}
\label{ch:data}

Kerangka di bab ini tidak datang dari satu kuliah. Ia muncul pertama kali di kuliah supervised
learning pada semester pertama, lalu kembali di kuliah unsupervised, di kuliah dasar-dasar
machine learning, dan di bagian pembuka kuliah geometric deep learning. Karena dipakai oleh
begitu banyak kuliah, saya menaruhnya di depan.

Semua yang ada di buku ini, dari LSTM sampai PINN, berdiri di atas satu gagasan sederhana:
kita punya contoh, dan kita ingin aturan yang juga berlaku untuk contoh yang belum pernah
dilihat. Bab ini membahas gagasan itu secara jujur. Apa yang sebenarnya kita minimalkan?
Mengapa model yang sempurna di data latih bisa gagal di data baru? Dan mengapa dimensi tinggi
membuat semuanya lebih sulit?

\section{Tiga cara belajar}

Bayangkan tiga murid. Murid pertama belajar dengan kunci jawaban: setiap soal latihan punya
jawaban benar, dan ia menyesuaikan cara berpikirnya sampai jawabannya cocok. Murid kedua hanya
diberi setumpuk soal tanpa jawaban. Ia mengelompokkan soal yang mirip dan menemukan polanya
sendiri. Murid ketiga tidak diberi soal sama sekali. Ia mencoba sesuatu, lalu gurunya hanya
berkata ``bagus'' atau ``kurang'', kadang baru setelah banyak langkah.

Murid pertama mewakili \istilah{supervised learning}, di mana setiap input $\vx$ datang bersama
target $y$. Murid kedua mewakili \istilah{unsupervised learning}, di mana hanya ada $\vx$ dan kita
mencari strukturnya: kelompok, arah utama, atau bentuk distribusinya. Murid ketiga mewakili
\istilah{reinforcement learning}, tempat sebuah agen bertindak, menerima \istilah{reward}, dan harus
menebak tindakan mana yang pantas dipuji (\cref{ch:rl}). Di antara ketiganya ada bentuk
campuran. \istilah{Self-supervised learning}, misalnya, membuat target dari data itu sendiri, seperti
menebak kata yang ditutup dalam sebuah kalimat. Hampir semua model bahasa besar dilatih dengan
cara ini. Bab ini berfokus pada murid pertama, karena kerangka berpikirnya dipakai di mana-mana.

\section{Kerangka supervised learning}

Kita menganggap data berasal dari sebuah distribusi gabungan $p(\vx,y)$ yang tidak kita ketahui.
Yang kita miliki hanyalah sampel $\mathcal{D}=\{(\vx_n,y_n)\}_{n=1}^N$ yang diambil secara
independen dari distribusi itu, asumsi yang biasa disebut \istilah{i.i.d.} Kita lalu memilih sebuah
model $f(\vx;\vw)$ dari suatu \istilah{model class} $\mathcal{F}$, dan sebuah \istilah{loss function}
$L(y,f(\vx))$ yang mengukur seberapa buruk tebakan kita. Dua besaran kemudian menjadi pusat
perhatian:
\begin{align}
R(f) &= \E_{(\vx,y)\sim p}\big[L(y,f(\vx))\big] = \int L(y,f(\vx))\,p(\vx,y)\dd\vx\dd y, \label{eq:risk}\\
\hat R_N(f) &= \frac1N\sum_{n=1}^N L(y_n,f(\vx_n)). \label{eq:emprisk}
\end{align}
Yang pertama disebut \istilah{generalization error} atau risiko sejati. Yang kedua disebut
\istilah{empirical risk} atau \istilah{training error}. Yang sebenarnya kita inginkan adalah $R(f)$
kecil, tetapi yang bisa kita hitung hanya $\hat R_N(f)$. \istilah{Empirical risk minimization}
(ERM) memilih $\hat f=\argmin_{f\in\mathcal{F}}\hat R_N(f)$ dan berharap $R(\hat f)$ ikut kecil.
Hampir seluruh teori generalisasi pada dasarnya bertanya kapan harapan ini masuk akal.

Seandainya kita mengetahui $p(\vx,y)$ dengan sempurna, ada prediktor terbaik yang mungkin, yaitu
\istilah{Bayes-optimal predictor} $f^*$, dengan risiko $R^*=R(f^*)$ yang disebut \istilah{Bayes risk}.
Untuk loss kuadrat, prediktor itu adalah rata-rata bersyarat $f^*(\vx)=\E[y\mid\vx]$. Alasannya
mudah dilihat. Tetapkan satu $\vx$ dan cari bilangan $c$ yang meminimalkan $\E[(y-c)^2\mid\vx]$.
Turunan terhadap $c$ adalah $-2\,\E[y-c\mid\vx]$, yang bernilai nol tepat ketika $c=\E[y\mid\vx]$.
Kalau loss diganti dengan $|y-c|$, jawaban optimalnya berubah menjadi median bersyarat. Jadi
pilihan loss menentukan statistik apa yang akan ditiru oleh model. Untuk klasifikasi dengan
loss 0--1, prediktor Bayes memilih kelas dengan peluang bersyarat terbesar. Bayes risk umumnya
tidak nol, karena data sendiri berisik: dua pasien dengan hasil laboratorium yang sama bisa saja
punya diagnosis yang berbeda.

\section{Dari asumsi noise ke fungsi loss}

Kuliah dasar-dasar machine learning menekankan satu hal yang sering terlewat: fungsi loss bukan
pilihan sembarang. Ia lahir dari asumsi tentang bagaimana data dibangkitkan. Misalkan
$y=f(\vx;\vw)+\varepsilon$ dengan noise $\varepsilon\sim\N(0,\sigma^2)$. Likelihood satu titik data
adalah
\[
p(y_n\mid\vx_n,\vw)=\frac{1}{\sqrt{2\pi\sigma^2}}\exp\!\Big(-\frac{(y_n-f(\vx_n;\vw))^2}{2\sigma^2}\Big).
\]
Karena data i.i.d., likelihood seluruh data adalah hasil kali, dan negatif log-likelihoodnya
menjadi
\begin{equation}
-\log p(\mathcal{D}\mid\vw)=\frac{1}{2\sigma^2}\sum_{n=1}^N\big(y_n-f(\vx_n;\vw)\big)^2+\frac N2\log(2\pi\sigma^2).
\end{equation}
Suku kedua tidak bergantung pada $\vw$. Maka \istilah{maximum likelihood} dengan noise Gauss sama
persis dengan meminimalkan \istilah{mean squared error}. Dengan jalan yang sama, noise Laplace
menghasilkan loss mutlak $|y-f|$, keluaran Bernoulli dengan peluang $\sigma(f)$ menghasilkan
\istilah{binary cross-entropy}, dan keluaran kategorikal dengan softmax menghasilkan
\istilah{cross-entropy} (\cref{tab:noiseloss}).

Kasus Bernoulli layak diturunkan juga, karena loss ini dipakai di hampir semua pengklasifikasi.
Misalkan label $y\in\{0,1\}$ dan model mengeluarkan peluang $p=\sigma(f(\vx;\vw))$ bahwa $y=1$. Peluang
satu label bisa ditulis dalam satu rumus, $p(y\mid\vx)=p^{\,y}(1-p)^{1-y}$, karena untuk $y=1$ rumus itu
menjadi $p$ dan untuk $y=0$ menjadi $1-p$. Negatif log-likelihood seluruh data adalah
\begin{equation}
-\log p(\mathcal D\mid\vw)=-\sum_{n=1}^N\Big[y_n\log p_n+(1-y_n)\log(1-p_n)\Big],
\end{equation}
dan itulah binary cross-entropy. Sebagai contoh, kalau model memberi peluang $0{,}9$ untuk contoh yang
labelnya satu, loss-nya $-\ln0{,}9\approx0{,}105$. Kalau peluangnya hanya $0{,}1$, loss-nya
$-\ln0{,}1\approx2{,}30$, lebih dari dua puluh kali lebih besar. Loss ini menghukum keyakinan yang
salah dengan keras sekali, dan sifat itulah yang membuat model terdorong memberi peluang yang
terkalibrasi.

\begin{table}[ht]
\centering\small
\caption{Asumsi noise dan loss yang dihasilkannya.}
\label{tab:noiseloss}
\begin{tabular}{@{}lll@{}}
\toprule
Asumsi pada $y\mid\vx$ & Loss dari $-\log$ likelihood & Sifat \\
\midrule
Gauss $\N(f,\sigma^2)$ & kuadrat, $(y-f)^2$ & peka terhadap pencilan \\
Laplace & mutlak, $|y-f|$ & lebih tahan pencilan \\
Bernoulli, $p=\sigma(f)$ & binary cross-entropy & klasifikasi dua kelas \\
Kategorikal, $p=\mathrm{softmax}(f)$ & cross-entropy & klasifikasi banyak kelas \\
\bottomrule
\end{tabular}
\end{table}

Timbangan dapur memberi gambaran yang pas. Kalau timbangan Anda kadang meleset sedikit ke atas
atau ke bawah secara acak, rata-rata dari beberapa kali menimbang adalah tebakan yang baik, dan
itulah loss kuadrat. Tetapi kalau sesekali seekor kucing melompat ke atas timbangan, satu hasil
bisa meleset jauh sekali. Rata-rata akan ikut tertarik oleh angka itu, sedangkan median tidak.
Noise berekor tebal seperti Laplace menghasilkan loss mutlak, yang tidak terlalu peduli pada
satu pencilan.

\section{Prior sebagai regularisasi}

Maximum likelihood hanya mendengarkan data. Kalau datanya sedikit, model mudah percaya pada
kebetulan. Pendekatan Bayesian menambahkan \istilah{prior} $p(\vw)$, yaitu keyakinan awal tentang
parameter sebelum melihat data. Teorema Bayes kemudian memberi \istilah{posterior}
\[
p(\vw\mid\mathcal{D})=\frac{p(\mathcal{D}\mid\vw)\,p(\vw)}{p(\mathcal{D})},
\]
dan estimasi \istilah{maximum a posteriori} (MAP) memaksimalkannya. Dalam bentuk logaritma,
\begin{equation}
\hat\vw_{\mathrm{MAP}}=\argmin_{\vw}\Big[-\log p(\mathcal{D}\mid\vw)-\log p(\vw)\Big].
\end{equation}
Dengan prior Gauss $\vw\sim\N(\bm 0,\alpha^{-1}\bm I)$, suku kedua menjadi $\tfrac\alpha2\|\vw\|_2^2$.
Itulah \istilah{weight decay}, atau \istilah{ridge regression} dalam bahasa statistik. Dengan prior
Laplace, suku kedua menjadi $\lambda\|\vw\|_1$, yang dikenal sebagai lasso \citep{tibshirani1996}.
Lasso punya sifat yang menarik: ia cenderung membuat banyak bobot bernilai tepat nol. Sifat ini
nanti menjadi inti metode SINDy untuk menemukan persamaan dari data (\cref{ch:penemuan}).

Dalam simulasi, kita sering tahu banyak hal sebelum melihat data. Aliran harus kekal massa,
temperatur dalam kelvin tidak boleh negatif, dan solusi harus simetris kalau geometrinya
simetris. Semua itu bisa dibaca sebagai prior yang kuat. PINN (\cref{ch:pinn}) memasukkan
persamaan diferensial ke dalam loss dengan cara yang mirip suku $-\log p(\vw)$ di atas, hanya
saja ``prior''-nya menyangkut fungsi keluaran, bukan bobot.

\section{Mengukur generalisasi}

Risiko sejati \eqref{eq:risk} tidak bisa dihitung, jadi kita memperkirakannya. Cara paling
sederhana adalah menyisihkan sebagian data yang tidak pernah disentuh selama pelatihan, lalu
menghitung loss rata-rata di sana. Karena titik-titik uji itu independen dari model yang sudah
jadi, rata-ratanya adalah estimator tak bias untuk $R(\hat f)$.

Kesulitan muncul karena kita juga harus memilih \istilah{hyperparameter}: jumlah lapisan, laju
belajar, kekuatan regularisasi. Kalau data uji dipakai untuk memilih, ia ikut bocor ke dalam
model dan estimasinya menjadi terlalu optimistis. Karena itu data biasanya dibagi tiga: data
latih, data validasi untuk memilih hyperparameter, dan data uji yang disentuh sekali saja di
akhir. Kalau data sedikit, \istilah{$k$-fold cross-validation} membagi data menjadi $k$ bagian dan
melatih model $k$ kali, setiap kali dengan satu bagian sebagai validasi. Rata-rata skornya
menjadi estimasi generalisasi, dengan biaya $k$ kali pelatihan. Perlu diingat bahwa
cross-validation menghasilkan estimasi kinerja, bukan satu model. Model akhir biasanya dilatih
ulang pada seluruh data latih.

Ada beberapa kebocoran yang sering terjadi tanpa disadari. Menormalisasi seluruh data, termasuk
data uji, sebelum dibagi berarti statistik data uji ikut masuk ke pelatihan. Rata-rata dan
simpangan baku seharusnya dihitung dari data latih saja. Untuk deret waktu atau hasil simulasi
yang berurutan, membagi data secara acak juga menyesatkan. Titik-titik yang berdekatan dalam
waktu hampir identik, sehingga ujian menjadi terlalu mudah. Data seperti ini sebaiknya dibagi
menurut waktu atau menurut kasus simulasi.

\section{Bias, varians, dan dua bentuk kurva}

Ambil loss kuadrat dan anggap $y=g(\vx)+\varepsilon$ dengan $\E[\varepsilon]=0$ dan
$\mathrm{Var}[\varepsilon]=\sigma^2$. Model $\hat f_{\mathcal D}$ bergantung pada data latih $\mathcal D$
yang acak. Untuk satu titik $\vx$, ekspektasi galatnya terhadap semua kemungkinan data latih dan
noise terurai menjadi tiga suku \citep{geman1992}:
\begin{equation}
\E\big[(y-\hat f_{\mathcal D}(\vx))^2\big]
=\underbrace{\big(g(\vx)-\E_{\mathcal D}[\hat f_{\mathcal D}(\vx)]\big)^2}_{\text{bias}^2}
+\underbrace{\E_{\mathcal D}\big[(\hat f_{\mathcal D}(\vx)-\E_{\mathcal D}[\hat f_{\mathcal D}(\vx)])^2\big]}_{\text{varians}}
+\underbrace{\sigma^2}_{\text{noise}}.
\end{equation}
Bias mengukur seberapa jauh rata-rata model dari kebenaran. Varians mengukur seberapa banyak
model berubah kalau data latihnya diganti. Noise tidak bisa dikurangi oleh model apa pun.

Penurunannya hanya memakai satu trik: tambahkan dan kurangkan besaran yang sama. Tulis singkat
$\hat f=\hat f_{\mathcal D}(\vx)$, $\bar f=\E_{\mathcal D}[\hat f]$, dan $g=g(\vx)$. Karena
$y=g+\varepsilon$,
\[
y-\hat f=\varepsilon+(g-\bar f)+(\bar f-\hat f).
\]
Kuadratkan, lalu ambil ekspektasi. Noise $\varepsilon$ berasal dari titik uji dan independen dari data
latih, sehingga semua suku silang yang memuat $\varepsilon$ hilang karena $\E[\varepsilon]=0$. Suku silang
$2(g-\bar f)\E_{\mathcal D}[\bar f-\hat f]$ juga hilang, karena $g-\bar f$ tidak acak dan
$\E_{\mathcal D}[\hat f]=\bar f$. Yang tersisa adalah tiga kuadrat: $\E[\varepsilon^2]=\sigma^2$,
$(g-\bar f)^2$, dan $\E_{\mathcal D}[(\hat f-\bar f)^2]$, persis tiga suku di atas.

Contoh angka membuatnya konkret. Misalkan nilai sebenarnya di sebuah titik $g=2$, noise bervarians
$\sigma^2=0{,}25$, dan kita melatih model yang sama pada tiga himpunan data latih berbeda, yang memberi
prediksi $1{,}6$, $2{,}0$, dan $2{,}1$ di titik itu. Rata-rata prediksinya $\bar f=1{,}9$, sehingga bias
kuadratnya $(2-1{,}9)^2=0{,}01$. Variansnya
$\tfrac13\big[(1{,}6-1{,}9)^2+(2{,}0-1{,}9)^2+(2{,}1-1{,}9)^2\big]=\tfrac13(0{,}09+0{,}01+0{,}04)\approx0{,}047$.
Galat kuadrat harapannya $0{,}01+0{,}047+0{,}25\approx0{,}31$, dan di sini noise adalah bagian terbesar.
Model yang lebih rumit bisa saja menurunkan bias, tetapi tidak akan menyentuh $0{,}25$ itu.

Papan dart membantu membayangkannya. Pemain dengan bias kecil dan varians kecil mengumpulkan
semua anak panah di tengah. Pemain dengan bias besar tetapi varians kecil mengumpulkannya rapat
di pojok kiri atas: konsisten, tetapi salah arah. Pemain dengan bias kecil tetapi varians besar
menyebar ke mana-mana, meski rata-ratanya tepat di tengah. Model yang terlalu sederhana mirip
pemain kedua. Model yang terlalu rumit mirip pemain ketiga, karena setiap set data latih
membuatnya melempar ke tempat yang berbeda.

Gambaran klasik mengatakan bahwa ketika kompleksitas model bertambah, bias turun dan varians
naik, sehingga galat uji membentuk huruf U dengan titik terbaik di tengah. Gambaran ini benar
untuk banyak model klasik, tetapi jaringan saraf modern sering punya parameter jauh lebih banyak
daripada jumlah data, dan tetap generalisasi dengan baik. \citet{belkin2019} menunjukkan bahwa
setelah \istilah{interpolation threshold}, yaitu titik ketika model baru saja mampu menghafal semua
data latih, galat uji bisa turun lagi. Kurva lengkapnya disebut \istilah{double descent}
(\cref{fig:doubledescent}). Sebelumnya, \citet{zhang2017} memperlihatkan bahwa jaringan yang sama
bisa menghafal label acak dengan sempurna. Artinya, kapasitas mentah tidak menjelaskan
generalisasi. Yang menentukan adalah solusi mana yang dipilih oleh algoritme pelatihan di antara
banyak solusi yang sama-sama menghafal data. Stochastic gradient descent tampaknya cenderung
memilih solusi yang halus, sebuah topik yang muncul lagi di \cref{ch:optimasi}.

\begin{figure}[ht]
\centering
\begin{tikzpicture}
\begin{axis}[width=0.48\linewidth,height=5cm,axis lines=left,xtick=\empty,ytick=\empty,
  xlabel={\small kompleksitas model},ylabel={\small galat},ymin=0,ymax=1.1,xmin=0,xmax=10,
  title={\small klasik},clip=false]
  \addplot[biru,thick,domain=0.3:10,samples=60] {0.9*exp(-0.5*x)+0.05};
  \addplot[oranye,thick,domain=0.3:10,samples=60] {0.9*exp(-0.5*x)+0.05+0.006*x^2};
  \node[biru,font=\scriptsize] at (axis cs:8.3,0.12) {latih};
  \node[oranye,font=\scriptsize] at (axis cs:8.3,0.62) {uji};
\end{axis}
\end{tikzpicture}\hfill
\begin{tikzpicture}
\begin{axis}[width=0.48\linewidth,height=5cm,axis lines=left,xtick={5},xticklabels={\scriptsize ambang interpolasi},ytick=\empty,
  xlabel={\small jumlah parameter},ymin=0,ymax=1.1,xmin=0,xmax=10,title={\small double descent},clip=false]
  \addplot[biru,thick,domain=0.3:10,samples=80] {max(0.9*exp(-0.6*x)-0.04*max(x-4,0),0.01)};
  \addplot[oranye,thick,domain=0.3:4.6,samples=60] {0.5*exp(-0.5*x)+0.15+0.08/(5.05-x)};
  \addplot[oranye,thick,domain=5.4:10,samples=60] {0.12+0.08/(x-4.95)};
  \draw[dashed,abu] (axis cs:5,0) -- (axis cs:5,1.05);
\end{axis}
\end{tikzpicture}
\caption{Kiri, kurva U klasik. Kanan, double descent: setelah model mampu menghafal semua data
latih, galat uji bisa turun lagi. Kedua kurva bersifat skematik.}
\label{fig:doubledescent}
\end{figure}

Kuliah geometric deep learning dan teori pembelajaran statistik memakai cara lain untuk
memecah galat, yang lebih umum. Misalkan $f^*$ prediktor Bayes, $f_{\mathcal F}$ prediktor terbaik
di dalam model class, $\hat f$ hasil ERM, dan $\tilde f$ hasil algoritme optimasi yang sungguh
kita jalankan. Maka
\begin{equation}
R(\tilde f)-R^*=\underbrace{R(f_{\mathcal F})-R^*}_{\text{aproksimasi}}
+\underbrace{R(\hat f)-R(f_{\mathcal F})}_{\text{estimasi}}
+\underbrace{R(\tilde f)-R(\hat f)}_{\text{optimasi}}.
\end{equation}
Model class yang lebih besar mengecilkan galat aproksimasi, tetapi bisa membesarkan galat
estimasi, karena semakin banyak fungsi yang kebetulan cocok dengan data. Galat optimasi muncul
karena kita hampir tidak pernah benar-benar menemukan minimum global.

\section{Kutukan dimensi}

Ada satu hal yang membuat machine learning di dimensi tinggi berbeda secara mendasar dari
intuisi kita di dimensi rendah. Namanya \istilah{curse of dimensionality}. Ambil kubus satuan
$[0,1]^d$ dan kubus kecil di tengahnya dengan sisi $0{,}9$. Fraksi volume di dalam kubus kecil
adalah $0{,}9^d$. Untuk $d=1$ hasilnya $0{,}9$. Untuk $d=10$ hasilnya sekitar $0{,}35$. Untuk
$d=100$ hasilnya kira-kira $2{,}7\times10^{-5}$. Hampir seluruh volume kubus berdimensi seratus
berada di kulit tipis dekat permukaannya. Titik acak hampir selalu dekat ke tepi, dan hampir
semua pasangan titik punya jarak yang mirip.

Akibatnya, metode lokal seperti $k$-nearest neighbour kehilangan makna. Untuk menutup ruang
berdimensi $d$ dengan ketelitian $\epsilon$, kita butuh sekitar $\epsilon^{-d}$ titik. Kuliah
geometric deep learning menyatakannya secara formal: untuk kelas fungsi Lipschitz di $d$ dimensi,
galat estimasi terbaik turun hanya seperti $N^{-1/d}$. Dengan $d=100$, menggandakan jumlah data
hampir tidak membantu.

Mencari kunci yang jatuh di lorong panjang masih mudah. Di lantai sebuah aula, butuh waktu jauh
lebih lama. Di gedung bertingkat yang penuh ruangan, kita mulai putus asa. Sekarang bayangkan
gedung dengan seratus arah yang berbeda: sebanyak apa pun senter yang dibawa, sebagian besar
ruang tetap gelap.

Lalu mengapa deep learning bisa bekerja pada gambar berukuran jutaan piksel? Ada dua jawaban
yang akan muncul berkali-kali di buku ini. Jawaban pertama adalah \istilah{manifold hypothesis}:
data nyata tidak memenuhi seluruh ruang. Gambar wajah manusia menempati semacam permukaan
berdimensi jauh lebih rendah di dalam ruang piksel. Kalau model bisa menemukan permukaan itu,
dimensi efektif masalahnya kecil. Jawaban kedua adalah simetri. Kita tahu bahwa kucing tetap
kucing kalau digeser sedikit ke kiri. Kalau arsitektur dibuat invarian atau ekuivarian terhadap
simetri seperti ini, ruang hipotesis menyusut drastis tanpa membuang solusi yang benar. Itulah
gagasan pokok \citet{bronstein2021}, yang kita bahas di \cref{ch:gdl}.

\section{Tidak ada makan siang gratis}

Teorema \istilah{no free lunch} \citep{wolpert1996} menyatakan bahwa kalau dirata-ratakan atas semua
kemungkinan masalah, tidak ada algoritme belajar yang lebih baik dari yang lain. Kedengarannya
pesimistis, tetapi pesannya praktis: keberhasilan selalu datang dari asumsi yang cocok dengan
masalahnya. Konvolusi bekerja untuk gambar karena gambar punya struktur lokal dan simetri geser.
PINN bekerja untuk masalah yang persamaannya kita ketahui. Tanpa asumsi, tidak ada
generalisasi.

\section{Mengukur kinerja klasifikasi}

Untuk klasifikasi dua kelas, hasil prediksi disusun dalam \istilah{confusion matrix} yang berisi
\istilah{true positive} (TP), \istilah{false positive} (FP), \istilah{true negative} (TN), dan
\istilah{false negative} (FN). Dari keempatnya diperoleh
\[
\text{precision}=\frac{TP}{TP+FP},\qquad
\text{recall}=\frac{TP}{TP+FN},\qquad
\text{specificity}=\frac{TN}{TN+FP}.
\]
Misalkan sebuah tes menandai 50 dari 1000 orang sebagai positif, dan 40 di antaranya memang sakit,
sementara total orang yang sakit ada 60. Maka $TP=40$, $FP=10$, dan $FN=20$, sehingga precision
$40/50=0{,}80$ dan recall $40/60\approx0{,}67$. Precision menjawab ``kalau tes bilang positif, seberapa
bisa dipercaya?'', recall menjawab ``dari semua yang sakit, berapa yang tertangkap?''. Keduanya sering
dirangkum dengan rata-rata harmonik, $F_1=2\cdot\frac{0{,}80\cdot0{,}67}{0{,}80+0{,}67}\approx0{,}73$.
Rata-rata harmonik dipakai karena ia tertarik ke nilai yang lebih kecil: model tidak bisa menutupi
recall yang buruk dengan precision yang sempurna.

Akurasi bisa menipu kalau kelasnya tidak seimbang. Pada data dengan 99 persen kelas negatif,
model yang selalu menjawab ``negatif'' punya akurasi 99 persen dan recall nol. Kurva ROC
memplot recall terhadap $1-\text{specificity}$ untuk semua ambang keputusan. Luas di bawahnya,
AUC, sama dengan peluang model memberi skor lebih tinggi pada satu contoh positif acak dibanding
satu contoh negatif acak.

\section{Generalisasi di luar kelas}

Pertanyaan-pertanyaan di bab ini akan muncul lagi di setiap bidang yang dibahas buku ini, dengan
wajah yang berbeda-beda.

Yang pertama adalah \istilah{distribution shift}: data saat model dipakai tidak berasal dari distribusi
yang sama dengan data latih. Model citra medis yang dilatih di satu rumah sakit bisa gagal di rumah
sakit lain dengan pemindai berbeda (\cref{ch:medis}). Model pengenal aktivitas yang dilatih pada
segelintir orang bisa gagal pada orang dengan gaya gerak lain (\cref{ch:pervasif}). Model pengganti
simulasi yang dilatih pada bilangan Reynolds antara 100 dan 1000 belum tentu benar pada 5000
(\cref{ch:operator}).

Yang kedua adalah kejujuran pembagian data. Di kimia, molekul dari satu seri senyawa hampir identik,
sehingga pembagian acak membuat model tampak jauh lebih baik dari kenyataannya (\cref{ch:hayati}). Di
sensor yang dikenakan, jendela-jendela dari orang yang sama berkorelasi kuat. Di simulasi, potret
berurutan dari satu kasus hampir sama. Pola umumnya: bagilah data menurut satuan yang benar-benar baru
saat model dipakai, entah itu pasien, orang, kerangka molekul, atau kasus simulasi.

Yang ketiga adalah mahalnya data. Satu label ahli radiologi, satu pengukuran assay, atau satu simulasi
CFD penuh bisa memakan waktu dan biaya besar. Di sinilah prior, dalam bentuk simetri, pengetahuan
domain, atau hukum fisika, menjadi berharga: ia mengurangi data yang dibutuhkan dan menjaga model
tetap masuk akal di luar daerah yang pernah dilihatnya.

\sumberbab{Bab ini merangkai materi dari Machine Learning: Basic Techniques, Supervised dan
Unsupervised Techniques, Statistical Learning Theory, serta bagian pembuka Deep Learning:
Geometric Techniques. Bacaan pendamping yang baik adalah \citet{bishop2006}, \citet{hastie2009},
dan \citet{bach2024}.}
